% !TeX root = ../main.tex


\chapter{高中数学实验教学资源的开发}
\section{开发的整体框架}
在正式开发实验教学资源之前，为确保开发与应用的合理性与可行性，要对整个过程进行思路规划，依据思路规划图，指导整个设计过程，见图 \ref{fig:开发框架思路图}。通过相关文献与媒体信息，进行开发需求分析。了解需求后着手资源设计，选取 “平面向量基本定理”、“三角函数图象平移伸缩变换”、“圆锥曲线的形成” 三个高中数学实验教学资源为内容，选择开发平台、搭建环境，开发资源、应用、反思与修改，最后发布使用。

\begin{figure}[htbp]
  \centering
  \includegraphics[width=0.8\textwidth]{开发框架思路图}
  \caption{开发框架思路图}
  \label{fig:开发框架思路图}
\end{figure}

\section{开发的设计原则}
在建构主义学习理论和认知同化理论的指导下，教学资源的使用对象是学生和教师，最终目的是深化学生的概念理解并促进知识迁移应用。为了实现上述目标，高中数学实验教学资源设计应当遵循以下设计原则：

\subsection{科学性原则}
开发高中数学实验教学资源的根基是科学性。在借助数字化软件演示实验的过程中，必须严格遵循数学的基本原理和逻辑规则。无论是构建数学模型还是展示数学现象，都要有坚实的数学理论支撑。在实验数据的呈现和分析上，也要保证科学严谨，杜绝任何虚假数据或错误结论，让学生在科学的环境中学习数学，培养正确的数学思维和科学态度，真实感受到信息科技的魅力，进而激发自身的学习探知欲。

\subsection{适切性原则}
首要任务是深度剖析教材内容以及学生的实际学情。精准把握高中数学教材各章节知识的重点、难点与易错点，同时充分考量不同年龄段学生的认知水平、 兴趣偏好以及已有的数学知识储备。基于此，借助 GeoGebra 和希沃白板设计出与教学内容紧密契合、符合学生能力层级的实验，确保实验教学资源能切实满足相应学生的学习需求，助力教学目标的达成。

\subsection{可视化原则}

数学知识往往抽象难懂，而 GeoGebra 和希沃白板强大的功能恰能有效化解这一难题。在开发实验资源过程中，要充分发挥它们的可视化优势。GeoGebra 能够将抽象的数学概念，如空间几何图形的旋转、向量的运算等，以生动直观的动态形式呈现出来。希沃白板则可通过图片、动画、视频等多媒体元素，把数学问题的解决过程可视化。比如在解析几何中，利用希沃白板展示圆锥曲线的形成过程，从平面截圆锥的动态演示到最终曲线方程的推导，让学生清晰地看到抽象数学知识的具象化呈现，从而降低理解难度，增强学生对数学知识的感性认识，提升学习效果。


\subsection{简洁性原则}
尽管 GeoGebra 和希沃白板功能丰富，但在开发实验教学资源时，应秉持简洁明了的原则。实验设计要避免过于复杂的操作流程和冗余的信息展示，确保学生能够迅速理解实验目的与操作方法。在实验教学资源设计中，教学素材的选取要精炼，以最直接有效的方式传达数学知识，让学生能专注于核心内容的学习，提高学习效率，避免因过度复杂的设计而分散学生的注意力。

\subsection{互动性原则}
高中数学教学强调学生的主动参与和自主探究。借助 GeoGebra 和希沃白板开发的实验教学资源，应具备良好的互动性。希沃白板的触摸交互功能，可让学生在课堂上直接参与数学问题的解决过程，亲身体验知识的生成。GeoGebra 互动式的实验教学，能够充分调动学生的学习积极性，培养学生的合作能力与创新思维，使学生真正成为学习的主体，在实验操作中实现思维的发散与体验的共享。


\section{开发的前期准备}
\subsection{开发平台}
本研究开发选择 GeoGebra 与希沃白板作为高中数学实验教学平台，源于二者在功能适配性、教学交互性及实践可行性上的互补优势。

GeoGebra 作为动态数学软件，其核心价值在于精准满足数学实验的学科特性。其不仅具备几何画板的动态作图功能，还整合了代数运算、3D 建模、概率统计等模块。软件内置的几何构造、代数运算及 3D 建模功能，能够动态呈现函数图象变换（如三角函数图象相位移动）、立体几何截面生成等传统课堂难以具象化的内容。例如，在讲解三角函数图象平移伸缩变换时，学生可通过调控参数滑动条，直观观察图象的变化过程。此外，GeoGebra 支持脚本指令与自定义工具，能快速构建复杂数学模型（如分形、微积分曲线），满足高中实验教学对深度探究的需求。其开源属性既降低学校软件采购成本，又支持教师根据教学需求二次开发特定实验模块。希沃白板的优势在于重构课堂教学场景。其交互组件与多屏协作功能，能将 GeoGebra 生成的数学模型快速转化为课堂互动素材。然而，现有的希沃白板交互实验演示资源并未开发涉及到圆锥曲线的形成探究实验。鉴于希沃白板的使用普及，本研究补充了圆锥曲线的形成探究实验的资源开发。二者的组合有效弥合了数学实验的技术门槛与教学落地之间的断层，帮助教师快速完成从数学建模到课堂实施的转换。


\subsection{环境搭建}

\begin{enumerate}
  \item 访问 GeoGebra 官网(\url{https://www.geogebra.org/})，在 “下载” 栏目中选择历史版本 4.0。Windows 用户选择.exe 安装包，Mac 用户下载.dmg 镜像文件。
  \item 运行下载的安装程序：Windows 用户双击执行文件后按向导完成路径配置； Mac 用户将应用图标拖入 “应用程序” 文件夹。
  \item 启动软件验证功能完整性：新建空白绘图区，输入基础函数(如 \(y = x^2\) )，若正常生成抛物线且工具栏交互无延迟，则表明安装成功。若界面显示异常，需检查系统显卡驱动是否支持 OpenGL 2.0 以上版本。
  \item 注意 4.0 版本兼容 32/64 位系统，但需预留至少 512MB 内存。教育版用户建议勾选 “CAS 计算器” 与 “3D 绘图” 模块以激活完整功能。
\end{enumerate}


\section{高中数学实验教学资源的开发示例}
\subsection{实验1——平面向量基本定理}

\begin{description}
  \item[实验准备]
  \hspace*{2em} 打开 GeoGebra 4.0（推荐电脑端操作），新建空白文件，选择 “代数区” 和 “图形区” 双视图模式，在工具栏确保勾选 “坐标轴” 显示，方便观察坐标。
  \item[开发步骤] \mbox{} \vspace*{-1.8em}
    \begin{enumerate}
      \item 绘制向量：在指令框中输入点 \(O\)，即坐标原点。选择向量工具，见图 \ref{fig:向量工具}，构造向量 \(\overrightarrow{OA}\)、\(\overrightarrow{OB}\) 与 \(\overrightarrow{OC}\)。右击向量 \(\overrightarrow{OA}\)、\(\overrightarrow{OB}\) 与 \(\overrightarrow{OC}\)，分别重命名为 \(\overrightarrow{e_1}\)、\(\overrightarrow{e_2}\) 与 \(\overrightarrow{a}\)。
      % \begin{figure}[htbp]
      %   \centering
      %   \includegraphics[width=0.3\textwidth]{向量工具}
      %   \caption{向量工具}
      %   \label{fig:向量工具}
      % \end{figure}
      \begin{figure}[htbp]
        \centering
        \begin{subfigure}[b]{0.45\textwidth}
          \centering
          \includegraphics[height=10cm]{向量工具}
          \caption{向量工具}
          \label{fig:向量工具}
        \end{subfigure}
        \hfill
        \begin{subfigure}[b]{0.45\textwidth}
          \centering
          \includegraphics[height=10cm]{平行线工具}
          \caption{平行线工具}
          \label{fig:平行线工具}
        \end{subfigure}
      \caption{向量工具与平行线工具}
    \end{figure}
      \item 绘制平行线：选择平行线工具，见图 \ref{fig:平行线工具}，先点击向量 \(\overrightarrow{OA}\)，再点击点 \(C\)，形成经过点 \(C\) 的向量 \(\overrightarrow{OA}\) 平行线 \((f)\)。同理得经过点 \(C\) 的向量 \(\overrightarrow{OB}\) 平行线 \((g)\)。右击直线，可依据需要设计相应颜色、显示状态，见图 \ref{fig:绘制平行线后示意图}。
      % \begin{figure}[htbp]
      %   \centering
      %   \includegraphics[width=0.3\textwidth]{平行线工具}
      %   \caption{平行线工具}
      %   \label{fig:平行线工具}
      % \end{figure}
      
      \begin{figure}[htbp]
        \centering
        \includegraphics[width=0.6\textwidth]{绘制平行线后示意图}
        \caption{绘制平行线后示意图}
        \label{fig:绘制平行线后示意图}
      \end{figure}
      
      \item 绘制交点及分解向量：在指令框中输入 “交点”，选择 “交点(<几何对象 圆锥曲线|向量|点|数值|多边形>，<几何对象 圆锥曲线|向量|点|数值|多边形>)”，左侧输入 \(f\)，右侧输入指令 “直线 \((\langle\) 点 \(1 \rangle\)，\(\langle\) 点 \(1 \rangle\))”，左 \(\langle\) 点 \(1 \rangle\) 输入 \(O\)，右 \(\langle\) 点 \(1 \rangle\) 输入 \(B\)，确定即得到交点，命名为点 \(N\)。同理，得直线 \(OA\) 与 \(g\) 交点 \(M\)。利用向量工具绘制向量 \(\overrightarrow{OM}\)、\(\overrightarrow{ON}\)，分别命名为 \(b\)、\(a_1\)，部分个体右击选择 “显示标签” 以隐藏标签，同时在 “属性” 中自由调整显示方式，效果见图 \ref{fig:分解向量效果图}。
      \begin{figure}[htbp]
        \centering
        \includegraphics[width=0.6\textwidth]{分解向量效果图}
        \caption{分解向量效果图}
        \label{fig:分解向量效果图}
      \end{figure}
      
      \item 进行计算：找到相关向量，例如向量 \(\overrightarrow{OB}(\overrightarrow{e_2})\) 与 \(\overrightarrow{ON}(\overrightarrow{a_1})\)，在指令框中输入 “\(x(a_1)/x(e_2)\)”，点击确定，得到数值 \(c\)，即平面向量分解的有序实数对之一。
      \item 设置文本：在指令框中输入 “如果”，选择 “如果 (\(\langle\) 条件 \(\rangle\)，\(\langle\) 结果 \(\rangle\)，\(\langle\) 否则 \(\rangle\))”，“条件” 输入 “\(c > 0\)”，“结果” 输入 “\(+\)” \(c\)，“否则” 输入 “\(c\)”，确定即产生文本 text1。创建文本，输入 “平面向量基本定理交互演示”，右击在 “属性” 中设置颜色、字体等；同理创建文本 “手动拖动 \(C\) 点位置，观察代数变化”。
      \item 显示分解内容：创建文本，勾选 “LaTeX 公式”，选择 “重音标记”，输入 \(a\)；输入 “=”，再引入空白格式框，输入 “\(x(b)/x(e_1)\)”；选择 “重音标记”，输入 \(\lambda e_1\)；点击 “对象”，选择 “text1”；选择 “重音标记”，输入 \(e_2\)，最终编辑见图 \ref{fig:向量线性分解文本编辑}；确定即产生文本 text2。
      \begin{figure}[htbp]
        \centering
        \includegraphics[width=0.5\textwidth]{向量线性分解文本编辑}
        \caption{向量线性分解文本编辑}
        \label{fig:向量线性分解文本编辑}
      \end{figure}
      
      \item 设置按钮：创建分布滑动条 \(n\)，区间设置为 “整数最小 \(1\)，最大 \(7\)，增量 \(1\)”；设置按钮，其一赋值 “\((n,n - 1)\)”，即后退按钮，其二赋值 “\((n,n + 1)\)”，即前进按钮；右击按钮，在 “属性” 中选择 “样式” 设置合适表现形式。
      \item 分图层显示操作步骤：分别选择平行线 \(f\)、\(g\)，在 “属性” - “高级” 的 “显示条件” 中输入 “\(n > 1\)”、“\(n > 2\)”；分别选择向量 \(\overrightarrow{OM}\)、\(\overrightarrow{ON}\)，在 “属性” - “高级” 的 “显示条件” 中输入 “\(n > 3\)”、“\(n > 4\)”；分别选择点 \(M\)、\(N\)，在 “属性” - “高级” 的 “显示条件” 中输入 “\(n > 3\)”、“\(n > 4\)”；选择 text2，在 “属性” - “高级” 的 “显示条件” 中输入 “\(n > 5\)”。
      \item 完善演示页面：创建文本，输入 “平面向量基本定理交互演示”，确定即产生文本 \(\text{text}3\)，右击在 “属性” 中设置颜色、字体等；同理创建文本 “手动拖动 \(C\) 点位置，观察代数变化”，即产生文本 \(\text{text}4\)。
      \item 实验预演：点击后退按钮恢复至初始状态，逐步点击前进按钮，观察界面演示至显示完整。拖动点 \(C\)，观察向量分解情况与代数变化。最终效果见图 \ref{fig:平面向量基本定理交互实验演示界面}。
      \begin{figure}[htbp]
        \centering
        \includegraphics[width=0.8\textwidth]{平面向量基本定理交互实验演示界面}
        \caption{平面向量基本定理交互实验演示界面}
        \label{fig:平面向量基本定理交互实验演示界面}
      \end{figure}
      
    \end{enumerate}
\end{description}


\subsection{实验2——三角函数图象平移伸缩变换}

\begin{description}
  \item[实验准备]
  \hspace*{2em} 打开 GeoGebra 4.0（推荐电脑端操作），新建空白文件，选择 “代数区” 和 “图形区” 双视图模式，在工具栏确保勾选 “网格” 显示，方便观察坐标。
  \item[开发步骤] \mbox{} \vspace*{-1.8em}
    \begin{enumerate}
      \item 绘制函数图象，同时展示平移变换、周期变换、振幅变换
        \begin{enumerate}
          \item 设置绘图区背景：右击空白处，选择绘图区，见图 \ref{fig:绘图区指示栏}，对 \(x\) 轴设置刻度间距为 \(\frac{\pi}{2}\)，见图 \ref{fig:设置x轴刻度间距}。
          \begin{figure}[htbp]
            \centering
            \includegraphics[width=0.2\textwidth]{绘图区指示栏}
            \caption{绘图区指示栏}
            \label{fig:绘图区指示栏}
          \end{figure}
          \begin{figure}[htbp]
            \centering
            \includegraphics[width=0.8\textwidth]{设置x轴刻度间距}
            \caption{设置$x$轴刻度间距}
            \label{fig:设置x轴刻度间距}
          \end{figure}
          
          \item 在输入框设置 \(A\)、\(\omega\)、\(\varphi\) 滑动条，设置适当最值与增量（本实验设置最值为 -5 与 5，增量为 0.1）。
          \item 输入基础函数：\(y = A\sin(\omega x + \varphi)\)，设置 “启动” 按钮，在按钮设置中输入脚本，见图 \ref{fig:启动按钮}。注意编辑脚本时，符号切换为半角输入，下面脚本操作类似，不再强调。点击 “启动” 按钮，函数 \(y = A\sin(\omega x + \varphi)\) 图象会随着 \(A\)、\(\omega\)、\(\varphi\) 滑动条的滑动而变化。
          \begin{figure}[htbp]
            \centering
            \includegraphics[width=0.5\textwidth]{启动按钮}
            \caption{启动按钮脚本}
            \label{fig:启动按钮}
          \end{figure}
          
          \item 设置 “暂停” 按钮，在按钮设置中输入脚本，见图 \ref{fig:暂停按钮脚本}。点击 “暂停” 按钮，图象会停止变化。
          \begin{figure}[htbp]
            \centering
            \includegraphics[width=0.5\textwidth]{暂停按钮脚本}
            \caption{暂停按钮脚本}
            \label{fig:暂停按钮脚本}
          \end{figure}
          
          \item 设置 “重置” 按钮，在按钮设置中输入脚本，见图 \ref{fig:重置按钮脚本}。点击 “重置” 按钮，图象会恢复初始状态。
          \begin{figure}[htbp]
            \centering
            \includegraphics[width=0.5\textwidth]{重置按钮脚本}
            \caption{重置按钮脚本}
            \label{fig:重置按钮脚本}
          \end{figure}
          
          \item 设置页数滑动条，将页面所有设置显示条件为“页数=2”，见图 \ref{fig:页数显示条件}。
          \begin{figure}[htbp]
            \centering
            \includegraphics[width=0.8\textwidth]{页数显示条件}
            \caption{页数显示条件}
            \label{fig:页数显示条件}
          \end{figure}
          
          \item 注意设置按钮脚本时，括号内左侧对象是相应滑动条名称，当滑动条命名不一样时，输入也不一样，下面操作类似，不再重复说明。最终效果见图 \ref{fig:三变量三角函数演示界面}。
          \begin{figure}[htbp]
            \centering
            \includegraphics[width=0.8\textwidth]{三变量三角函数演示界面}
            \caption{三变量三角函数演示界面}
            \label{fig:三变量三角函数演示界面}
          \end{figure}
          
        \end{enumerate}
      \item 设置增加页面和翻页脚本
        \begin{enumerate}
          \item 增加页数滑动条，以此增加显示页面。
          \item 设置 “上一页”、“下一页” 按钮，对 “上一页” 按钮设置脚本 “赋值(页数, 页数 - 1)”，对 “下一页” 按钮设置脚本 “赋值(页数, 页数 + 1)”。
        \end{enumerate}
      \item 平移变换演示
        \begin{enumerate}
          \item 创建相位参数：在代数区输入：\(\varphi = 0\) (滑动条设置范围 \(-2\pi\) 到 \(2\pi\))。
          \item 输入基础函数：\(y = \sin x\)，输入函数为：\(y = \sin(x + \varphi)\)。
          \item 创建 “向右平移 \(\frac{\pi}{3}\) 个单位长度”、“向左平移 \(\frac{\pi}{2}\) 个单位长度” 按钮，向右按钮设置脚本 “赋值 \((\varphi, \varphi - \frac{\pi}{3})\)”，向左按钮设置脚本 “赋值 \((\varphi, \varphi - \frac{\pi}{2})\)”。
          \item 创建 “重置” 按钮，设置脚本如图 \ref{fig:重置按钮脚本2}。
          \begin{figure}[htbp]
            \centering
            \includegraphics[width=0.5\textwidth]{重置按钮脚本 2}
            \caption{重置按钮脚本}
            \label{fig:重置按钮脚本2}
          \end{figure}
          
          \item 点击向右按钮，观察图象向右平移 \(\frac{\pi}{3}\) 个单位长度；点击向左按钮，观察图象向左平移 \(\frac{\pi}{2}\) 个单位长度；点击重置按钮，图象恢复初始状态。最终效果见图 \ref{fig:平移变换演示界面}。
          \begin{figure}[htbp]
            \centering
            \includegraphics[width=0.8\textwidth]{平移变换演示界面}
            \caption{平移变换演示界面}
            \label{fig:平移变换演示界面}
          \end{figure}
          
        \end{enumerate}
      \item 周期变换演示
        操作同平移变换演示，将参数由 \(\varphi\) 设置为 \(\omega\) 即可，同时设置按钮，相应改变按钮脚本，最终效果见图 \ref{fig:周期变换演示界面}。
        \begin{figure}[htbp]
          \centering
          \includegraphics[width=0.8\textwidth]{周期变换演示界面}
          \caption{周期变换演示界面}
          \label{fig:周期变换演示界面}
        \end{figure}
        
      \item 振幅变换演示
        操作同平移变换演示，将参数由 \(\varphi\) 设置为 \(A\) 即可，同时设置按钮，相应改变按钮脚本，最终效果见图 \ref{fig:振幅变换演示界面}。
        \begin{figure}[htbp]
          \centering
          \includegraphics[width=0.8\textwidth]{振幅变换演示界面}
          \caption{振幅变换演示界面}
          \label{fig:振幅变换演示界面}
        \end{figure}
        
      \item 综合变换演示
        操作同平移变换演示，将一个参数由 \(\varphi\) 设置为三个参数 \(A\)、\(\varphi\)、\(\omega\) 即可，同时设置按钮，相应改变按钮脚本，最终效果见图 \ref{fig:综合变换演示界面}。
        \begin{figure}[htbp]
          \centering
          \includegraphics[width=0.8\textwidth]{综合变换演示界面}
          \caption{综合变换演示界面}
          \label{fig:综合变换演示界面}
        \end{figure}
        
    \end{enumerate}
\end{description}


\subsection{实验3——圆锥曲线的形成}

\begin{description}
  \item[实验准备]
  \hspace*{2em} 打开网络画板软件（推荐电脑端操作），新建空白文件，选择 3D 视图，将画布颜色改为白色，将坐标面隐藏，显示坐标轴。
  \item[实验步骤] \mbox{} \vspace*{-1.8em}
    \begin{enumerate}
      \item 创建双圆锥：在 \(y\) 轴上建立一个自由点 \(A(14)\)，为下一步建立圆锥打下基础；使用圆台/圆锥/圆柱工具，并选中点 \(O\) 与点 \(A\)，建立一个圆锥(15)；测量 \(OA\) 之间的距离 \(m_0\) 作为圆锥的高(16)；建立变量尺 \(R(17)\) 作为圆锥的半径；再利用对称工具(18)，并选中建立的圆锥，建立一个以点 \(O\) 对称的新圆锥。见图 \ref{fig:双圆锥示意图}。
      \begin{figure}[htbp]
        \centering
        \includegraphics[width=0.8\textwidth]{双圆锥示意图}
        \caption{双圆锥示意图}
        \label{fig:双圆锥示意图}
      \end{figure}
      
      \item 创建截面：使用交点工具描点两个圆锥的交点(19)，利用对称工具找到上面圆锥底面的圆心(21)，连接圆心与点 \(A\)，建立一条线段(22)，在刚刚建立的线段上描一个动点(23)，利用动点以及 \(z\) 轴建立一个垂直于 \(z\) 轴的平面(24)。见图 \ref{fig:创建截面效果示意图}。
      \begin{figure}[htbp]
        \centering
        \includegraphics[width=0.8\textwidth]{创建截面效果示意图}
        \caption{创建截面效果示意图}
        \label{fig:创建截面效果示意图}
      \end{figure}
      
      \item 创建切割效果：建立一个变量 \(t(25)\)，利用刚建立的垂直平面以及变量 \(t(25)\)，建立一个旋转平面(26)；选中圆锥(15)以及旋转平面(26) 设置切割(28)；建立一个变量 \(d(29)\)，选择对称圆锥(18)建立一个切割效果(30)。见图 \ref{fig:切割效果操作界面图}。 
      \begin{figure}[htbp]
        \centering
        \includegraphics[width=\textwidth]{切割效果操作界面图}
        \caption{切割效果操作界面图}
        \label{fig:切割效果操作界面图}
      \end{figure}
      
      \item 显示截面边界线：建立一个参数动画平移(31)，以及圆锥与平面的两个路径贯穿线(32)(33)，利用平移将刚刚的路径进行平移(34)(35)，平移量为变量 \(c(36)\)。 点击右下角的动作按钮，建立一个动作圆(37)，变量为 \(t = 0\)。点击右下角的动作按钮，建立一个动作椭圆，变量为 \(t = 0.2\)。见图 \ref{fig:建立动作操作界面图}。
      \begin{figure}[htbp]
        \centering
        \includegraphics[width=0.8\textwidth]{建立动作操作界面图}
        \caption{建立动作操作界面图}
        \label{fig:建立动作操作界面图}
      \end{figure}
      
      \item 同理创建抛物线和双曲线的动作。
      \item 显示截面倾斜角：在贯穿线上描点 \(S(41)\)，连接 \(O\) 与 \(S(42)\)，建立 \(S\) 与(22)线段的垂线(43)。同时，建立线段(43)与线段(22)的交点(44)，测量角 \(\angle OTS(45)\)，并显示。见图 \ref{fig:截面倾斜角显示示意图}。
      \begin{figure}[htbp]
        \centering
        \includegraphics[width=0.6\textwidth]{截面倾斜角显示示意图}
        \caption{截面倾斜角显示示意图}
        \label{fig:截面倾斜角显示示意图}
      \end{figure}
      
      \item 创建按钮：建立重置与角度按钮，关联变量 \(t\)，并将需要显示的进行显示，不需要显示的进行隐藏。
      \item 创建文本：按需建立相应文字叙述的文本。
      \item 调整界面：最后对显示的按钮进行调整，使界面和谐。最终效果见图 \ref{fig:圆锥曲线形成的实验演示界面}。
      \begin{figure}[htbp]
        \centering
        \includegraphics[width=\textwidth]{圆锥曲线形成的实验演示界面}
        \caption{圆锥曲线形成的实验演示界面}
        \label{fig:圆锥曲线形成的实验演示界面}
      \end{figure}
      
    \end{enumerate}
\end{description}


\subsection{实验开发小结——基于GeoGebra软件}

本研究开发的三个实验资源中，有两个是由GeoGebra制作。在实际设计开发过程中，对GeoGebra软件的资源开发使用有独到的感悟理解，叙述如下：
\begin{enumerate}
  \item 紧扣教学目标  
  优先选择教材中抽象难懂的知识点（如函数变换、立体几何展开、圆锥曲线定义），通过动态演示将抽象概念具象化。

  \item 常用核心功能说明  
    \begin{enumerate}
      \item 参数联动：用多个滑动条控制复合函数（如设置分别控制二次函数的开口、对称轴、顶点）。
      \item 条件显示：设置图形随参数变化的显隐规则（如勾选“显示渐近线”时才出现双曲线渐近线）。
      \item 动态文本：让公式随参数实时变化（如显示当前滑动条值对应的函数解析式）。
      \item 多轨迹对比：同时追踪不同初始条件的动点轨迹（如改变初速度观察不同抛物线）。
      \item 轨迹分析工具：添加动态切线、曲率半径等分析工具（适合导数等实验）。
      \item 逆向构造：从轨迹反推构造过程（如已知椭圆轨迹，分解展示两焦点与动点关系）。
    \end{enumerate}

  \item 感悟之外  
  更多有关GeoGebra软件的功能说明可在官网介绍与自身动手实操中发现与学习。
\end{enumerate}


\subsection{实验资源共享操作}

本研究开发了两个GeoGebra实验资源和一个希沃白板实验资源。

针对GeoGebra实验资源，可以在GeoGebra官网注册账号，将文件保存到 GeoGebra 账户中，选择生成链接进行分享；被共享者可以直接进入“查看链接”实现对文件的预览和基础使用，也可登录账号进入“编辑链接”进行创意编辑，保存为己用的实验资源。若想实现大平台资源共享，可以自行选择按照一定分类发布文件到GeoGebra资源库。

针对希沃白板实验资源，鉴于登录账户可以体验交互功能，并展示动画，故建议下载希沃白板软件，同时也方便后续共享协作。开发完成的资源上传到云空间，通过点击共享链接进入，登录账号，实现对文件的预览和获取，后续进行二次创作。

以下为本研究开发实验资源的链接：

\begin{enumerate}[label=实验 \arabic*.]
  \item 平面向量基本定理
      \begin{itemize}
          \item 查看链接：\url{https://www.geogebra.org/m/txgwnw62}
          \item 编辑链接：\url{https://www.geogebra.org/material/edit/id/ef9kwqkjm}
      \end{itemize}
  \item 三角函数图象平移伸缩变换
      \begin{itemize}
          \item 查看链接：\url{https://www.geogebra.org/m/emteakpk}
          \item 编辑链接：\url{https://www.geogebra.org/material/edit/id/ete6grgwc}
      \end{itemize}
  \item 圆锥曲线的形成
      \begin{itemize}
          \item 预览/接收链接：\url{https://13.r302.cc/0AA0bGJ?platform=enpc&channel=copylink}
      \end{itemize}
\end{enumerate}